\section{Implementasi dan Simulasi}
Bab ini menguraikan implementasi \algname{} dalam aplikasi \textit{Command Line Interface} (CLI), hasil pengujian simulasi enkripsi dan dekripsi menggunakan berkas uji yang \textit{reproducible}, mekanisme penanganan \textit{padding}, serta verifikasi integritas data berbasis \textit{Encrypt-then-MAC}.

\subsection{Deskripsi Program}
Aplikasi CLI (\texttt{cipher\_cli}) dibangun menggunakan C++17 tanpa \textit{dependency} pustaka kriptografi pihak ketiga. Komponen inti (fungsi putaran $F$, ekspansi kunci, lima mode cipher, \textit{padding}, dan CMAC) diimplementasikan langsung pada kode sumber, dengan pemanfaatan OpenMP untuk mempercepat mode ECB, CTR, serta dekripsi CBC dan CFB.

\subsubsection{Antarmuka Perintah dan Format Kontainer}
Program menyediakan tiga perintah operasional:
\begin{itemize}
  \item \texttt{keygen [-f] -o <path>}: Membangkitkan kunci acak 128 bit (32 hex) menggunakan CSPRNG sistem operasi (\texttt{getentropy}) dengan izin berkas POSIX \texttt{0600} (\texttt{S\_IRUSR | S\_IWUSR}).
  \item \texttt{encrypt [-f] -m <mode> (-k <hex> | --key-file <path>) [--iv <hex>]}\linebreak
    \texttt{(-i <in> | -t <teks>) -o <out>}: Mengenkripsi berkas atau teks masukan ke berkas keluaran.
  \item \texttt{decrypt [-f] (-k <hex> | --key-file <path>) -i <in> -o <out>}: Mendekripsi berkas kembali ke data aslinya (mode dibaca otomatis dari \textit{header}).
\end{itemize}

Seluruh data terenkripsi dikemas dalam format biner \textbf{Container v1} dengan \textit{overhead} tetap 38 byte.

\begin{table}[H]
  \centering
  \caption{Format Kontainer v1}\label{tab:container_layout}
  \begin{tabular}{lcp{9.5cm}}
    \toprule
    Field & Ukuran & Deskripsi \\
    \midrule
    \textit{Magic Bytes} & 4 Byte & Pengenal format ASCII \texttt{"IF15"} (\texttt{0x49 0x46 0x31 0x35}) \\
    \textit{Version}     & 1 Byte & Nomor versi protokol (\texttt{0x01}) \\
    \textit{Mode ID}     & 1 Byte & ID mode: \texttt{0x01}=ECB, \texttt{0x02}=CBC, \texttt{0x03}=CFB, \texttt{0x04}=OFB, \texttt{0x05}=CTR \\
    \textit{IV}          & 16 Byte & Vektor inisialisasi / counter awal (diisi byte nol pada mode ECB) \\
    \textit{Payload}     & $N$ Byte & Data terenkripsi ($N$ kelipatan 16 pada ECB/CBC, $N = L$ pada mode aliran) \\
    \textit{CMAC Tag}    & 16 Byte & Tag autentikasi integritas RFC 4493 atas \textit{header} dan \textit{payload} \\
    \bottomrule
  \end{tabular}
\end{table}

Penyertaan \textit{header} (Mode ID dan IV) ke dalam tag CMAC mencegah \textit{Mode-Downgrade Attack} (penyerang memanipulasi byte mode ke ECB) dan \textit{IV Bit-Flipping Attack} pada mode CBC.

\subsubsection{Konvensi \textit{Exit Code} dan Keamanan Runtime}
Untuk memudahkan otomasi pengujian, \texttt{cipher\_cli} menetapkan \textit{exit code} terstruktur: \texttt{0} (sukses), \texttt{1} (kesalahan argumen/CLI), \texttt{2} (kesalahan I/O berkas), \texttt{3} (format rusak/\texttt{FormatError}), dan \texttt{4} (autentikasi gagal/\texttt{AuthError}).

\subsection{Hasil Uji Enkripsi dan Dekripsi}
Pengujian simulasi dilakukan menggunakan berkas uji yang telah tersedia di repositori (\texttt{test-results/samples/}) agar hasil bersifat deterministik dan \textit{reproducible}:
\begin{itemize}
  \item \textbf{Berkas Teks}: \texttt{teks.txt} (1.116 byte, SHA-256: \texttt{0118caf99b62bdf9\dots{}}), memuat teks multibaris, karakter non-ASCII, dan panjang bukan kelipatan 16.
  \item \textbf{Berkas Biner}: \texttt{gambar.png} (949 byte, SHA-256: \texttt{5acafa9454163dc6\dots{}}), berkas citra format PNG.
\end{itemize}

Pengujian menggunakan pasangan kunci dan IV tetap:
\begin{align*}
  K_{\text{master}} &= \texttt{000102030405060708090a0b0c0d0e0f} \\
  \text{IV} &= \texttt{0f0e0d0c0b0a09080706050403020100} \quad \text{(khusus non-ECB)}
\end{align*}

Prosedur eksekusi enkripsi, dekripsi, dan verifikasi kesamaan biner pada terminal dapat direproduksi melalui perintah:
\begin{lstlisting}[language=bash]
# Enkripsi dan dekripsi berkas teks pada mode CBC
./cipher_cli encrypt -m cbc -k 000102030405060708090a0b0c0d0e0f --iv 0f0e0d0c0b0a09080706050403020100 -i test-results/samples/teks.txt -o out.enc
./cipher_cli decrypt -k 000102030405060708090a0b0c0d0e0f -i out.enc -o out.dec

# Verifikasi kesamaan bit-demi-bit
cmp -s test-results/samples/teks.txt out.dec && echo "Identik 100%"
\end{lstlisting}

\subsubsection{Berkas Teks}
Pada berkas \texttt{teks.txt} (1.116 byte), mode ECB dan CBC menerapkan \textit{padding} PKCS\#7 sehingga \textit{payload} menjadi 1.120 byte dan kontainer berukuran 1.158 byte ($22 + 1.120 + 16$). Pada mode CFB, OFB, dan CTR, \textit{payload} tepat 1.116 byte dan kontainer berukuran 1.154 byte ($22 + 1.116 + 16$). Seluruh hasil dekripsi pada kelima mode terverifikasi 100\% identik dengan berkas sumber via perintah \texttt{cmp}.

\subsubsection{Berkas Biner}
Pada berkas \texttt{gambar.png} (949 byte), mode ECB dan CBC menghasilkan \textit{payload} 960 byte dan kontainer 998 byte. Mode aliran menghasilkan \textit{payload} 949 byte dan kontainer 987 byte. Seluruh hasil dekripsi memiliki nilai hash SHA-256 yang identik dan gambar PNG dapat dibuka secara visual tanpa kerusakan struktur \textit{chunk}.

Rangkuman kuantitatif hasil pengujian \textit{round-trip} disajikan pada Tabel~\ref{tab:roundtrip_results}.

\begin{table}[H]
  \centering
  \caption{Hasil Uji Round-Trip Enkripsi dan Dekripsi Berkas Uji}\label{tab:roundtrip_results}
  \begin{tabular}{llcccc}
    \toprule
    Berkas & Mode & Plainteks & Kontainer & Payload & Status Dekripsi \\
    \midrule
    \texttt{teks.txt}   & ECB & 1.116 B & 1.158 B & 1.120 B & Identik (\texttt{pass}) \\
    \texttt{teks.txt}   & CBC & 1.116 B & 1.158 B & 1.120 B & Identik (\texttt{pass}) \\
    \texttt{teks.txt}   & CFB & 1.116 B & 1.154 B & 1.116 B & Identik (\texttt{pass}) \\
    \texttt{teks.txt}   & OFB & 1.116 B & 1.154 B & 1.116 B & Identik (\texttt{pass}) \\
    \texttt{teks.txt}   & CTR & 1.116 B & 1.154 B & 1.116 B & Identik (\texttt{pass}) \\
    \midrule
    \texttt{gambar.png} & ECB & 949 B   & 998 B   & 960 B   & Identik (\texttt{pass}) \\
    \texttt{gambar.png} & CBC & 949 B   & 998 B   & 960 B   & Identik (\texttt{pass}) \\
    \texttt{gambar.png} & CFB & 949 B   & 987 B   & 949 B   & Identik (\texttt{pass}) \\
    \texttt{gambar.png} & OFB & 949 B   & 987 B   & 949 B   & Identik (\texttt{pass}) \\
    \texttt{gambar.png} & CTR & 949 B   & 987 B   & 949 B   & Identik (\texttt{pass}) \\
    \bottomrule
  \end{tabular}
\end{table}

\subsection{Penanganan Padding}
Mode berbasis blok (ECB dan CBC) mewajibkan data merupakan kelipatan dari ukuran blok (16 byte). Program mengimplementasikan skema \textit{padding} PKCS\#7 sesuai RFC 5652~\cite{rfc5652}. Untuk panjang data $L$ byte, jumlah byte \textit{padding} $p$ bernilai:
\begin{equation}
  p = 16 - (L \bmod 16)
\end{equation}
di mana setiap byte \textit{padding} bernilai heksadesimal $p$:
\begin{itemize}
  \item \texttt{teks.txt}: $L = 1.116 \implies 1.116 \bmod 16 = 12 \implies p = 4$. Ditambahkan 4 byte \texttt{0x04 0x04 0x04 0x04}.
  \item \texttt{gambar.png}: $L = 949 \implies 949 \bmod 16 = 5 \implies p = 11$. Ditambahkan 11 byte \texttt{0x0B}.
  \item Data kelipatan 16 ($L \bmod 16 = 0$): Ditambahkan 1 blok penuh (16 byte \texttt{0x10}) untuk menjamin determinisme saat \textit{unpadding}.
\end{itemize}

Sebaliknya, mode CFB, OFB, dan CTR beroperasi sebagai \textit{stream cipher} yang melakukan XOR langsung antara \textit{keystream} dan data, sehingga tidak memerlukan \textit{padding} ($p = 0$).

\subsection{Verifikasi Integritas Data}
Integritas data dijamin menggunakan metode \textbf{Encrypt-then-MAC (EtM)} berbasis CMAC (RFC 4493)~\cite{rfc4493} dengan tag 128 bit (16 byte).

\subsubsection{Pemisahan Kunci dan Perbandingan Waktu-Konstan}
Sistem menerapkan dua prinsip perancangan kriptografi defensif:
\begin{enumerate}
  \item \textbf{Pemisahan Kunci (\textit{Key Separation})}: Kunci enkripsi ($K_{\text{enc}}$) dan kunci MAC ($K_{\text{mac}}$) diturunkan secara terpisah dari $K_{\text{master}}$ menggunakan fungsi pembangkit semu (PRF):
  \begin{equation}
    K_{\text{enc}} = \text{CIPH}_{K_{\text{master}}}(0^{15} \,\|\, \text{0x01}), \quad
    K_{\text{mac}} = \text{CIPH}_{K_{\text{master}}}(0^{15} \,\|\, \text{0x02})
  \end{equation}
  Pemisahan ini menjamin independensi matematis antara kerahasiaan dan integritas data.
  \item \textbf{Komparasi Waktu-Konstan}: Pengecekan tag CMAC menggunakan fungsi akumulasi bitwise OR branchless (\texttt{util::constant\_time\_equal}):
  \begin{equation}
    \Delta = \bigvee_{i=0}^{15} (T_{\text{dihitung}}[i] \oplus T_{\text{diterima}}[i])
  \end{equation}
  Fungsi ini membutuhkan waktu eksekusi yang konstan tanpa keluar lebih awal (\textit{early exit}), sehingga meniadakan serangan kanal samping waktu (\textit{timing attack}).
\end{enumerate}

\subsubsection{Prinsip Fail-Stop}
Pada proses dekripsi, tag CMAC diverifikasi \textbf{sebelum} mendekripsi \textit{ciphertext} atau memeriksa \textit{padding}. Jika verifikasi gagal, sistem langsung berhenti (\textit{fail-stop}). Mekanisme ini mengeliminasi peluang terjadinya serangan \textit{Padding Oracle} (Vaudenay) karena dekriptor tidak pernah memproses data yang terkorupsi.

